\uuid{vFm3}
\titre{Zéros de fonction}
\auteur{bodin ; hanani ; mzari}
\organisation{exo7}
\licence{CC BY-NC-SA 4.0}
\datecreate{2020-01-20}
\niveau{L1}
\module{Analyse}
\chapitre{Continuité, limite et étude de fonctions réelles}
\sousChapitre{Continuité : théorie}
\difficulte{1}
\theme{Théorèmes des valeurs intermédiaires, zéros de fonction}
\exo7qcmid{lille-1.continuite.020}

\contenu{
\begin{question}
\qclassification{123.01,123.02}
\qtags{motcle=zéros de fonction}
Soit $f(x)=x^2-7$. On applique la méthode de dichotomie sur l'intervalle $[2 ; 3]$.
On calcule $f(2,125)=-1,9375$ ; $f(2,5) = -0,75$ ; $f(2,625) = -0,109375$ ; $f(2,75) = 0,5625$. Quelles sont les assertions vraies ?
\begin{answers}
    \bad{$f$ s'annule sur $[2 ; 2,5]$ et sur $[2,5 ; 3]$.}

    \good{$f$ s'annule sur $[2,5 ; 3]$.}

    \bad{$f$ s'annule sur $[2,75 ; 3]$.}

    \good{$2,6 \le \sqrt{7} \le 2,8$}

\end{answers}
\begin{explanations}
La fonctions $f$ est continue, strictement croissante et s'annule en $\sqrt{7}$.
Comme $f(2,625) < 0$ et $f(2,75) > 0$, alors
$2,625 < \sqrt{7} < 2,75$.
\end{explanations}
\end{question}
}
